\section{Часть 2. Оценка защищённости сторонних программных систем}
\label{sec:part2}

В качестве объекта оценки выступает приложение Secure Notes Lab (результат части~1): Flask, SQLite, локальный HTTP на порту 5050, модули регистрации/авторизации, заметок и секретов.

Разработанная методика включает девять этапов:

\begin{enumerate}
	\item инвентаризация компонентов (ОС, runtime, СУБД, библиотеки, протоколы);
	\item построение модели угроз (активы, субъекты, типовые атаки);
	\item анализ ОС и окружения (права на файлы, пользователь процесса);
	\item анализ компонентов и взаимодействий (аутентификация, авторизация, CSRF, шифрование);
	\item анализ сетевой политики (интерфейс прослушивания, TLS, cookie);
	\item анализ конфиденциальности данных (классификация, at-rest encryption, изоляция);
	\item анализ организационных мер (администрирование, бэкапы, ротация ключей);
	\item практическая проверка сценариев использования и негативных тестов;
	\item формирование выводов, мер улучшения и правил безопасного использования.
\end{enumerate}

Уровни рисков: высокий / средний / низкий~\mdash  по возможности эксплуатации и ущербу для конфиденциальности, целостности и доступности. Полный текст методики~\mdash  файл \texttt{lab2/p2/methodology.md}.

Результаты применения методики согласуются с анализом части~1 и дополнительно фиксируют высокий риск при компрометации файлов \texttt{instance/} (ключ + БД), средний риск брутфорса пароля без rate limiting, средний риск MitM при выносе сервиса в сеть без TLS и низкий риск SQLi/XSS при текущей реализации (ORM + autoescape), при условии отсутствия отключения экранирования.

Комплекс мер по усовершенствованию:

\begin{enumerate}
	\item вывод ключей шифрования из каталога данных; per-user DEK на базе KDF от пароля;
	\item ограничение частоты попыток входа и блокировка учётной записи;
	\item публикация только через reverse-proxy с TLS и \texttt{Secure}-cookie;
	\item журнал аудита обращений к секретам;
	\item заголовки безопасности (CSP, X-Content-Type-Options);
	\item регулярные обновления зависимостей и ротация секретов приложения.
\end{enumerate}

Рекомендации по безопасному использованию:

\begin{itemize}
	\item не публиковать порт 5050 в Интернет без прокси и TLS;
	\item не включать каталог \texttt{instance/} в публичные репозитории;
	\item использовать длинные уникальные пароли;
	\item запускать от отдельного непривилегированного пользователя ОС;
	\item после демонстрации удалять или ротировать учебные ключи и БД.
\end{itemize}
